\chapter{Qualité de marché et régulation}\label{ch:mx:market-quality-and-regulation}

Vingt ans d'automatisation ont resserré les fourchettes et multiplié les messages, et les régulateurs ne s'accordent toujours pas sur le point de savoir si les marchés y gagnent.
Les données sont abondantes~; ce qu'elles mesurent, et ce qu'elles ne peuvent pas mesurer, fait l'objet de ce chapitre. Il définit les mesures, lit les principales études sur
le \omterm{def:mx:market-quality-and-regulation:hft}{trading algorithmique} et le \omterm{def:mx:market-quality-and-regulation:hft}{trading à haute fréquence}, le flash crash et les courses à la latence, puis évalue une règle introduite à dessein sur un panel simulé de jours-titres, où
l'effet vrai de la règle est connu, pour montrer ce que retrouve une estimation en \omterm{def:mx:tick-size-queues-and-priority:did}{doubles différences} et ce qu'une comparaison avant-après fausse.

\section{Mesurer la liquidité et la qualité de marché}

\begin{definition}[Qualité de marché]\label{def:mx:market-quality-and-regulation:quality}
La \emph{qualité de marché}\index{qualité de marché} (market quality) est l'ensemble des propriétés par lesquelles un marché sert ceux qui y négocient~: le coût de la négociation (\omterm{def:mx:decomposing-the-spread:quoted}{fourchette affichée},
fourchettes effectives et réalisées, impact sur les prix), la quantité disponible (profondeur, résilience) et le caractère informatif des prix (à quel point ils restent proches
d'une marche aléatoire, à quelle vitesse ils intègrent l'information).
\end{definition}

Les mesures sont celles des chapitres et livres précédents~: la \omterm{def:mx:decomposing-the-spread:quoted}{fourchette affichée} (chapitre~5), les fourchettes effectives et réalisées et leur écart,
l'impact sur les prix (le livre~1, chapitre~10), la profondeur et la résilience (chapitre~3), le ratio d'illiquidité d'Amihud (le livre~7, chapitre~7) et le ratio
de variance, soit la variance des rendements sur un long intervalle rapportée à la somme des variances sur les intervalles courts qui le composent (égal à un pour une marche aléatoire, inférieur à un
lorsque les cours rebondissent entre le cours acheteur et le cours vendeur). \texttt{firm\_mktquality.day\_measures} calcule ces mesures pour un jour-titre à partir du chemin du meilleur niveau du carnet d'ordres, des transactions et
du nombre de messages d'ordres.

\begin{definition}[Ratio ordres sur transactions]\label{def:mx:market-quality-and-regulation:otr}
Le \emph{ratio ordres sur transactions}\index{ratio ordres sur transactions} (order-to-trade ratio) d'un participant ou d'une action est le nombre de messages d'ordres (nouveaux ordres, modifications,
annulations) par transaction sur une période.
\end{definition}

\section{Trading algorithmique et à haute fréquence dans les études empiriques}

\begin{definition}[Trading algorithmique, trading à haute fréquence]\label{def:mx:market-quality-and-regulation:hft}
Le \emph{trading algorithmique}\index{trading algorithmique} (algorithmic trading) est une négociation dans laquelle un algorithme informatique détermine les paramètres des ordres (s'il faut les envoyer,
leur moment, leur prix ou leur quantité, et la façon de les gérer ensuite) avec peu ou pas d'intervention humaine. \emph{trading à haute
fréquence}\index{trading à haute fréquence} (high-frequency trading) est une forme de trading algorithmique qui repose sur une infrastructure minimisant la latence (colocalisation, hébergement de proximité ou
accès direct à haute vitesse), décide sans intervention humaine pour chaque ordre et envoie de nombreux ordres, cotations ou annulations dans la journée.
\end{definition}

Ces définitions suivent l'article~4(1)(39) et (40) de MiFID~II, qui excluent les systèmes qui se bornent à acheminer des ordres ou à les traiter sans en décider les
paramètres. Les données sont surtout favorables à la liquidité en temps normal. Hendershott, Jones et Menkveld (2011) ont exploité comme instrument l'automatisation de la diffusion des cotations de la New York Stock Exchange
en 2003, qui a accru le \omterm{def:mx:market-quality-and-regulation:hft}{trading algorithmique}~: pour les grandes actions, le \omterm{def:mx:market-quality-and-regulation:hft}{trading algorithmique} a resserré
les fourchettes, réduit la sélection adverse et rendu les cotations plus informatives. Brogaard, Hendershott et Riordan (2014) ont montré que les traders à haute fréquence
favorisent l'efficience des prix en négociant dans le sens des variations permanentes des cours et à contre-courant des erreurs de cotation transitoires, principalement par leurs
ordres consommateurs de liquidité, tandis que leurs ordres fournisseurs de liquidité sont soumis à la sélection adverse. Menkveld (2013) a décrit un grand trader à haute
fréquence~: il a perdu sur son inventaire, a encaissé la fourchette, a participé à 8.1~\% des transactions sur le marché historique et à 64.4~\% sur un nouvel entrant, et a été passif dans
quatre transactions sur cinq.

\section{La littérature sur le flash crash}

\begin{definition}[Trading de la patate chaude]\label{def:mx:market-quality-and-regulation:hotpotato}
Le \emph{trading de la patate chaude}\index{trading de la patate chaude} (hot potato) est le fait d'intermédiaires qui se renvoient rapidement entre eux les mêmes positions, ce qui
gonfle les volumes sans transférer le risque à quiconque souhaite le détenir.
\end{definition}

Le 6 mai 2010 (le livre~1, chapitre~31), un ensemble de fonds communs de placement a commencé à vendre 75\,000 contrats E-mini, soit environ 4.1~milliards de dollars, avec un algorithme réglé pour
négocier 9~\% du volume de la minute précédente, sans égard au prix ni au temps~: l'\omterm{def:mx:benchmark-algorithms:vwap}{algorithme de participation} du chapitre~16, avec la rétroaction que le chapitre
décrivait. Le rapport du personnel de la CFTC et de la SEC constate que, faute d'acheteurs fondamentaux, les traders à haute fréquence se sont mis à acheter et à revendre des contrats les uns aux
autres~; entre 2:45~:13 et 2:45~:27, ils ont échangé plus de 27\,000 contrats, soit environ 49~\% du volume, tout en n'achetant qu'environ 200 contrats nets~: un effet de volume de la patate chaude
que l'algorithme suivant le volume a pris pour de la liquidité. Kirilenko, Kyle, Samadi et Tuzun (2017), à partir de la piste d'audit de l'E-mini, ont constaté que les
intermédiaires non désignés les plus actifs n'ont pas modifié leur manière de négocier lorsque les cours ont baissé. Le volume n'est pas la liquidité~: une mesure de \omterm{def:mx:market-quality-and-regulation:quality}{qualité de marché} qui
compte les transactions aurait classé les pires minutes du flash crash comme les plus liquides.

\section{Courses à la latence et ralentisseurs dans les études empiriques}

Aquilina, Budish et O'Neill (2022) ont utilisé les données de messages des bourses, qui enregistrent les tentatives échouées de négociation ou d'annulation aussi bien que les réussies, pour observer
directement les courses d'arbitrage de latence~: pour les actions du FTSE~100, environ une course par minute et par symbole, une course modale durant de 5 à 10 millionièmes de seconde, et
environ un cinquième de l'ensemble des volumes négociés. Le chapitre~10 a montré que les \omterm{def:mx:auction-mechanics-and-theory:fba}{enchères par lots fréquentes} suppriment la course par construction~; les ralentisseurs et l'économie de la course
elle-même relèvent du livre~11. Pour ce chapitre, la leçon est méthodologique~: les courses étaient invisibles dans les données de carnet d'ordres que la plupart des études
utilisent, et une mesure ne peut voir que ce que ses données enregistrent.

\section{Règles sur les algorithmes et ratios ordres sur transactions}

Le concept release de la SEC de 2010 sur la structure des marchés d'actions demandait des commentaires sur le \omterm{def:mx:market-quality-and-regulation:hft}{trading à haute fréquence}, l'acheminement des ordres, les liens avec les données de marché et
la liquidité non affichée~; MiFID~II a suivi avec ses définitions et une exigence imposée aux plateformes de suivre, pour chaque membre, le
ratio ordres non exécutés sur transactions.

\begin{dated}{2026-09}{Ratios ordres sur transactions dans l'UE}\label{dat:mx:market-quality-and-regulation:otr}
La Commission Delegated Regulation (EU) 2017/566 (RTS~9 au titre de MiFID~II) oblige les plateformes de négociation à calculer le ratio ordres non exécutés sur transactions
de chaque membre et participant, pour chaque instrument financier négocié dans un système à carnet d'ordres continu, dirigé par les cotations ou hybride~; il s'applique à compter de la
date à laquelle MiFID~II s'est appliqué. Le texte cité est le règlement tel qu'adopté.
\end{dated}

Qu'une telle règle améliore le marché est une question empirique, et la simulation y répond là où la vérité est connue. Vingt actions
d'activité différente négocient pendant seize journées de dix minutes dans \texttt{firm.agentmkt}. À partir du jour~8, deux choses se produisent~: le marché de chaque action devient plus
volatil et ses fournisseurs de liquidité envoient un quart d'ordres en moins (un changement global du marché, le facteur de confusion)~; et une action sur deux, classées par activité,
est soumise à un cap sur le \omterm{def:mx:market-quality-and-regulation:otr}{ratio ordres sur transactions} qui conduit ses fournisseurs de liquidité à envoyer 30~\% d'ordres en moins et à annuler moins. Chaque jour-titre traité après l'événement est
aussi rejoué sans le cap, avec le même flux d'ordres, de sorte que l'effet de la règle est mesuré directement.

\begin{omfigure}
\begin{tikzpicture}
\begin{axis}[width=9.6cm,height=5.4cm,xlabel={\small jour par rapport à la règle},ylabel={\small fourchette affichée (pb)},
  tick label style={font=\footnotesize},grid=major,grid style={gray!25},xmin=-8.5,xmax=7.5,ymin=1,ymax=3.6,
  legend columns=2,legend style={font=\footnotesize,at={(0.5,-0.3)},anchor=north,draw=none,fill=none,
  /tikz/every even column/.append style={column sep=8pt}},table/col sep=comma]
\addplot[omProp,very thick,mark=*,mark size=1.4pt] table[x=day,y=treated]{figdata/microstructure/24-market-quality-and-regulation/daily.csv};
\addplot[omDef,very thick,mark=square*,mark size=1.4pt] table[x=day,y=control]{figdata/microstructure/24-market-quality-and-regulation/daily.csv};
\addplot[gray,dashed] coordinates {(-0.5,1) (-0.5,3.6)};
\legend{actions traitées,actions témoins}
\end{axis}
\end{tikzpicture}

\omcaption{La \omterm{def:mx:decomposing-the-spread:quoted}{fourchette affichée} moyenne des dix actions traitées et des dix actions témoins, jour par jour. Les deux augmentent à l'événement (le changement global du marché)~; les actions traitées
augmentent davantage (la règle). Données~: \texttt{mx\_mq.study}.}\label{fig:mx:market-quality-and-regulation:daily}
\end{omfigure}

La \omterm{def:mx:decomposing-the-spread:quoted}{fourchette affichée} des actions traitées a augmenté de 1.60~points de base (erreur type 0.08) entre l'avant et l'après~; l'estimation en \omterm{def:mx:tick-size-queues-and-priority:did}{doubles différences}
(méthode du chapitre~6, avec des effets fixes par action et par jour et des erreurs regroupées par action, \texttt{firm\_mktquality.did}) vaut 0.89 (0.09)~; la
valeur vraie est 0.87 (0.03). La comparaison avant-après attribue à la règle le changement global de 0.71~points de base
(le \cref{tab:mx:market-quality-and-regulation:did}).

\omcode{code/firm/mktquality/firm_mktquality.py}{61}{71}{L'estimation en doubles différences avec effets fixes à deux voies~: l'indicatrice de traitement après l'événement, les indicatrices par action et par jour, et les erreurs types regroupées par action.}\label{lst:mx:market-quality-and-regulation:did}

\begin{table}[ht]
\centering
\small
\begin{tabular}{@{}lrrr@{}}
\toprule
mesure & avant-après (e.t.) & \omterm{def:mx:tick-size-queues-and-priority:did}{doubles différences} (e.t.) & valeur vraie (e.t.) \\
\midrule
\omterm{def:mx:decomposing-the-spread:quoted}{fourchette affichée} (pb) & 1.60 (0.08) & 0.89 (0.09) & 0.87 (0.03) \\
fourchette effective (pb) & 2.74 (0.12) & 1.31 (0.14) & 1.36 (0.06) \\
fourchette réalisée (pb) & 1.81 (0.08) & 1.06 (0.10) & 1.15 (0.07) \\
impact sur les prix (pb) & 0.93 (0.09) & 0.25 (0.10) & 0.21 (0.07) \\
profondeur (actions) & $-138$ (8) & $-20$ (13) & $-31$ (2) \\
ratio de variance & $-0.19$ (0.02) & $-0.13$ (0.06) & $-0.14$ (0.03) \\
\omterm{def:mx:market-quality-and-regulation:otr}{ratio ordres sur transactions} & $-2.64$ (0.05) & $-1.18$ (0.07) & $-1.19$ (0.02) \\
\bottomrule
\end{tabular}
\caption{L'effet de la règle sur les actions traitées~: la variation avant-après, l'estimation en \omterm{def:mx:tick-size-queues-and-priority:did}{doubles différences} et l'effet vrai issu des
mêmes journées simulées sans la règle. Données~: \texttt{mx\_mq.study}.}\label{tab:mx:market-quality-and-regulation:did}
\end{table}

La règle a fait ce qu'elle devait faire (1.2 message de moins par transaction) et a dégradé le marché sur toutes les autres mesures~: fourchettes plus larges, profondeur moindre, davantage
de rebond bid-ask. L'étude d'événement (\cref{fig:mx:market-quality-and-regulation:event}) vérifie le dispositif~: avant l'événement, les fourchettes des actions traitées et celles des actions témoins
évoluent ensemble (les sept coefficients restent à moins de 0.07~point de base de zéro), puis après lui l'écart s'ouvre immédiatement et se maintient.

\begin{omfigure}
\begin{tikzpicture}
\begin{axis}[width=9.6cm,height=5.4cm,xlabel={\small jour par rapport à la règle},ylabel={\small traitées moins témoins (pb)},
  tick label style={font=\footnotesize},grid=major,grid style={gray!25},xmin=-8.5,xmax=7.5,ymin=-0.4,ymax=1.6,table/col sep=comma]
\addplot[omDef,very thick,mark=*,mark size=1.4pt,error bars/.cd,y dir=both,y explicit] table[x=day,y=coef,y error expr=1.96*\thisrow{se}]{figdata/microstructure/24-market-quality-and-regulation/event.csv};
\addplot[black,dashed,domain=-8.5:7.5] {0};
\end{axis}
\end{tikzpicture}

\omcaption{Étude d'événement de la \omterm{def:mx:decomposing-the-spread:quoted}{fourchette affichée}~: la différence traitées moins témoins chaque jour par rapport à la veille de la règle, avec effets fixes par action et par jour
et intervalles à 95~\% regroupés par action. Données~: \texttt{mx\_mq.study}.}\label{fig:mx:market-quality-and-regulation:event}
\end{omfigure}

\section{Tutoriel~: la règle a-t-elle amélioré le marché~?}\label{tut:mx:market-quality-and-regulation:tut}

\begin{tutorial}
\textbf{Objectif.} Construire un panel action-journée de mesures de \omterm{def:mx:market-quality-and-regulation:quality}{qualité de marché} dans le marché simulé, imposer une règle à la moitié des actions, et l'estimer par
\omterm{def:mx:tick-size-queues-and-priority:did}{doubles différences}, en la confrontant à la comparaison avant-après et à la vérité.
\textbf{État final~:} \cref{fig:mx:market-quality-and-regulation:daily,fig:mx:market-quality-and-regulation:event},
\cref{tab:mx:market-quality-and-regulation:did}.
\begin{tutsteps}
\item \textbf{Mesures.} \texttt{firm\_mktquality.day\_measures(top, trades, messages)}.
\item \textbf{Panel.} \texttt{mx\_mq.config(i, d, rule)}, \texttt{session\_day}, \texttt{panel()}.
\item \textbf{Estimation.} \texttt{did}, \texttt{before\_after}, \texttt{event\_study}~; \texttt{study()}~; dessiner avec \texttt{fig\_mq.py}.
\end{tutsteps}
\textbf{Pour aller plus loin.} Traiter les dix actions les moins actives au lieu d'une sur deux et constater la défaillance de l'hypothèse de tendances parallèles
(exercice~7)~; calculer les mesures sur les données MIDAS de la SEC pour un vrai changement de règle.
\end{tutorial}

\section{Construction~: qualité de marché}\label{bld:mx:market-quality-and-regulation:mktquality}

\begin{build}
\bfield{Rôle} Mesurer les marchés et les règles~: les mesures de panel et les estimateurs qui sous-tendent ce chapitre et les coupe-circuits du chapitre~25.
\bfield{Interface} \texttt{day\_measures(top, trades, messages, h, grid, px\_per\_ccy)}, \texttt{did(y, treated, post, unit, period)},
\texttt{before\_after(y, treated, post, unit)}, \texttt{event\_study(y, treated, period, unit, base)}.
\bfield{Règles} Fourchettes en points de base du milieu~; champs de transaction et de meilleurs cours de firm.tape~; erreurs regroupées par unité de traitement.
\bfield{Tests d'acceptation} \texttt{code/firm/mktquality/tests/}~: les mesures sur un marché à la main~; un effet implanté retrouvé par \omterm{def:mx:tick-size-queues-and-priority:did}{doubles différences}
et manqué par la comparaison avant-après~; une période préalable plate dans l'étude d'événement.
\bfield{Pour aller plus loin} Résilience~; contrôles synthétiques~; adoption échelonnée.
\end{build}

\begin{omsources}
\item personnel de la CFTC et de la SEC, ``Findings regarding the market events of May 6, 2010'', 2010.
\item SEC, Concept release on equity market structure, Release No.~34-61358, 2010.
\item T.~Hendershott, C.~M.~Jones et A.~J.~Menkveld, ``Does algorithmic trading improve liquidity?'', \emph{Journal of Finance} 66(1), 2011.
\item A.~J.~Menkveld, ``High frequency trading and the new market makers'', \emph{Journal of Financial Markets} 16(4), 2013.
\item J.~Brogaard, T.~Hendershott et R.~Riordan, ``High-frequency trading and price discovery'', \emph{Review of Financial Studies} 27(8), 2014.
\item A.~Kirilenko, A.~S.~Kyle, M.~Samadi et T.~Tuzun, ``The flash crash: high-frequency trading in an electronic market'', \emph{Journal of
Finance} 72(3), 2017.
\item M.~Aquilina, E.~Budish et P.~O'Neill, ``Quantifying the high-frequency trading `arms race''', \emph{Quarterly Journal of Economics} 137(1),
2022.
\item Directive 2014/65/EU (MiFID~II), article 4(1)(39)--(40); Commission Delegated Regulation (EU) 2017/566.
\end{omsources}

\section{Exercices}

\begin{exercise}[$\star$]\label{exo:mx:market-quality-and-regulation:1}
À partir des quatre moyennes par groupe de la \omterm{def:mx:decomposing-the-spread:quoted}{fourchette affichée} (traitées 1.58 avant et 3.18 après, témoins 1.56 et 2.27), calculer l'estimation en \omterm{def:mx:tick-size-queues-and-priority:did}{doubles différences}.
\end{exercise}

\begin{exercise}[$\star$]\label{exo:mx:market-quality-and-regulation:2}
Une action-journée compte 900 messages d'ordres et 250 transactions. Quel est son \omterm{def:mx:market-quality-and-regulation:otr}{ratio ordres sur transactions}, et que deviendrait-il avec l'effet de la règle, soit $-1.18$~?
\end{exercise}

\begin{exercise}[$\star$]\label{exo:mx:market-quality-and-regulation:3}
Pourquoi un ratio de variance inférieur à un est-il le signe d'un rebond bid-ask, et qu'a fait la règle à ce ratio~?
\end{exercise}

\begin{exercise}[$\star\star$]\label{exo:mx:market-quality-and-regulation:4}
Pourquoi les erreurs types sont-elles regroupées par action et non par action-journée~?
\end{exercise}

\begin{exercise}[$\star\star$]\label{exo:mx:market-quality-and-regulation:5}
Pourquoi le volume des traders à haute fréquence pendant le flash crash a-t-il induit en erreur un \omterm{def:mx:benchmark-algorithms:vwap}{algorithme de participation}~?
\end{exercise}

\begin{exercise}[$\star\star$]\label{exo:mx:market-quality-and-regulation:6}
Pourquoi les données du carnet d'ordres ne peuvent-elles pas montrer les courses à la latence, et quelles données le peuvent~?
\end{exercise}

\begin{exercise}[$\star\star\star$]\label{exo:mx:market-quality-and-regulation:7}
\emph{Programmation.} Traiter les dix actions les moins actives au lieu d'une sur deux. Qu'estiment les \omterm{def:mx:tick-size-queues-and-priority:did}{doubles différences} pour la \omterm{def:mx:decomposing-the-spread:quoted}{fourchette affichée}, et
pourquoi passent-elles à côté de la vérité~?
\end{exercise}

\begin{exercise}[$\star\star\star$]\label{exo:mx:market-quality-and-regulation:8}
\emph{Trouver la faille.} «~Depuis la règle, les fourchettes des actions qu'elle couvre se sont élargies de 1.6~point de base~: la règle a coûté 1.6~point de base aux investisseurs.~»
\end{exercise}

\section{Problème~: la règle a-t-elle amélioré le marché~?}

\begin{problem}[{Problème du week-end --- la règle a-t-elle amélioré le marché~?}]\label{pb:mx:market-quality-and-regulation:1}
Un régulateur a plafonné les \omterm{def:mx:market-quality-and-regulation:otr}{ratios ordres sur transactions} sur la moitié des actions d'un marché. Il est demandé de l'évaluer.

\medskip\noindent\textbf{Partie I --- Mesures.}
\begin{enumerate}
\item Définir la \omterm{def:mx:market-quality-and-regulation:quality}{qualité de marché} et énumérer ses mesures.
\item Comment les fourchettes effective et réalisée et l'impact sur les prix sont-ils liés~?
\item Définir le \omterm{def:mx:market-quality-and-regulation:otr}{ratio ordres sur transactions}, le \omterm{def:mx:market-quality-and-regulation:hft}{trading algorithmique} et le \omterm{def:mx:market-quality-and-regulation:hft}{trading à haute fréquence}.
\item Que mesure le ratio de variance~?
\end{enumerate}

\medskip\noindent\textbf{Partie II --- Les preuves.}
\begin{enumerate}[resume]
\item Qu'ont trouvé Hendershott, Jones et Menkveld, et comment l'ont-ils identifié~?
\item Qu'ont trouvé Brogaard, Hendershott et Riordan sur la \omterm{def:mx:fragmentation-and-routing:discovery}{formation des prix}~?
\item Décrire le programme de vente du flash crash et le volume du \omterm{def:mx:market-quality-and-regulation:hotpotato}{trading de la patate chaude}.
\item Qu'ont mesuré Aquilina, Budish et O'Neill, et pourquoi ont-ils pu le faire~?
\end{enumerate}

\medskip\noindent\textbf{Partie III --- La règle simulée.}
\begin{enumerate}[resume]
\item Décrire le panel, le facteur confondant et la règle.
\item Comment la vérité est-elle mesurée~?
\item Écrire la régression à double effets fixes.
\item \emph{Énoncer le résultat nommé}~: l'estimation par \omterm{def:mx:tick-size-queues-and-priority:did}{doubles différences} de l'effet de la règle sur les \omterm{def:mx:decomposing-the-spread:quoted}{fourchettes affichées}, avec son erreur type, face à
l'estimation avant-après qui ignore le groupe de contrôle.
\item Que montre l'étude d'événement avant l'événement, et pourquoi cela importe-t-il~?
\end{enumerate}

\medskip\noindent\textbf{Partie IV --- Jugement.}
\begin{enumerate}[resume]
\item La règle a-t-elle amélioré le marché~? Utiliser le tableau.
\item Que se passe-t-il si les actions traitées sont les moins actives~?
\item Quelle mesure a le plus bougé en valeur relative, et pourquoi~?
\item De quoi aurait-on besoin pour évaluer une telle règle sur des données réelles~?
\item Pourquoi le volume est-il une mauvaise mesure de la liquidité~?
\item Que dit l'encadré daté sur les \omterm{def:mx:market-quality-and-regulation:otr}{ratios ordres sur transactions} dans l'UE~?
\item En une phrase~: qu'apporte un groupe de contrôle à une évaluation~?
\end{enumerate}
\end{problem}

\section{Questions d'entretien}

\begin{interviewq}[$\star$ \iqroles{trader}]\label{iq:mx:market-quality-and-regulation:1}
Les fourchettes sur votre action ont doublé après un changement de règle. Comment savez-vous que c'est bien la règle qui l'a provoqué~?
\end{interviewq}

\begin{interviewq}[$\star\star$ \iqroles{chercheur}]\label{iq:mx:market-quality-and-regulation:2}
Expliquez les \omterm{def:mx:tick-size-queues-and-priority:did}{doubles différences} et leur hypothèse clé. Comment les testeriez-vous~?
\end{interviewq}

\begin{interviewq}[$\star\star$ \iqroles{chercheur}]\label{iq:mx:market-quality-and-regulation:3}
Le \omterm{def:mx:market-quality-and-regulation:hft}{trading à haute fréquence} améliore-t-il la \omterm{def:mx:market-quality-and-regulation:quality}{qualité de marché}~? Que disent les preuves, et que ne disent-elles pas~?
\end{interviewq}

\begin{interviewq}[$\star\star$ \iqroles{risque}]\label{iq:mx:market-quality-and-regulation:4}
Quelles leçons le flash crash offre-t-il à un \omterm{def:mx:benchmark-algorithms:algo}{algorithme d'exécution}~?
\end{interviewq}

\begin{interviewq}[$\star\star$ \iqroles{développeur}]\label{iq:mx:market-quality-and-regulation:5}
Votre société doit rester sous le \omterm{def:mx:market-quality-and-regulation:otr}{ratio ordres sur transactions} d'une plateforme de négociation. Comment le surveilleriez-vous et le feriez-vous respecter dans vos systèmes de trading~?
\end{interviewq}

\begin{interviewq}[$\star\star\star$ \iqroles{chercheur}]\label{iq:mx:market-quality-and-regulation:6}
Comment mesureriez-vous les courses à la latence sur une plateforme de négociation, et que vous faudrait-il obtenir de celle-ci~?
\end{interviewq}
